\documentclass[10pt,a4paper]{amsart}
\usepackage{amsmath,amssymb,amsthm,graphicx}
% Retained preamble for the Stacks Project chapter snapshot in this directory.
% The chapter's own preamble.tex is a copy of the shared Stacks Project preamble; the block
% below is sliced MECHANICALLY, line for line, from it. Nothing was added, removed or
% redefined; the wrapper class, the xr-hyper/hyperref wiring and the externaldocument list
% are omitted because the assembler supplies the document class.
%
% lines taken from preamble.tex: 3, 5, 9, 11, 12, 156, 157, 158, ... (51 lines total)
%
% package coverage restored from the chapter's own preamble (2026-10-13): xy, xy-2cell

\usepackage[all]{xy}
\xyoption{2cell}
\UseAllTwocells

\newenvironment{reference}{\comment}{\endcomment}
\newenvironment{slogan}{\comment}{\endcomment}
\newenvironment{history}{\comment}{\endcomment}
\theoremstyle{plain}
\newtheorem{theorem}[subsection]{Theorem}
\newtheorem{proposition}[subsection]{Proposition}
\newtheorem{lemma}[subsection]{Lemma}
\theoremstyle{definition}
\newtheorem{definition}[subsection]{Definition}
\newtheorem{example}[subsection]{Example}
\newtheorem{exercise}[subsection]{Exercise}
\newtheorem{situation}[subsection]{Situation}
\theoremstyle{remark}
\newtheorem{remark}[subsection]{Remark}
\newtheorem{remarks}[subsection]{Remarks}
\numberwithin{equation}{subsection}
\def\lim{\mathop{\mathrm{lim}}\nolimits}
\def\colim{\mathop{\mathrm{colim}}\nolimits}
\def\Spec{\mathop{\mathrm{Spec}}}
\def\Hom{\mathop{\mathrm{Hom}}\nolimits}
\def\Ext{\mathop{\mathrm{Ext}}\nolimits}
\def\SheafHom{\mathop{\mathcal{H}\!\mathit{om}}\nolimits}
\def\SheafExt{\mathop{\mathcal{E}\!\mathit{xt}}\nolimits}
\def\Sch{\mathit{Sch}}
\def\Mor{\mathop{\mathrm{Mor}}\nolimits}
\def\Ob{\mathop{\mathrm{Ob}}\nolimits}
\def\Sh{\mathop{\mathit{Sh}}\nolimits}
\def\NL{\mathop{N\!L}\nolimits}
\def\CH{\mathop{\mathrm{CH}}\nolimits}
\def\proetale{{pro\text{-}\acute{e}tale}}
\def\etale{{\acute{e}tale}}
\def\QCoh{\mathit{QCoh}}
\def\Ker{\mathop{\mathrm{Ker}}}
\def\Im{\mathop{\mathrm{Im}}}
\def\Coker{\mathop{\mathrm{Coker}}}
\def\Coim{\mathop{\mathrm{Coim}}}
\DeclareMathSymbol{\boxtimes}{\mathbin}{AMSa}{"02}
\def\QCohstack{\mathcal{QC}\!\mathit{oh}}
\def\Cohstack{\mathcal{C}\!\mathit{oh}}
\def\Spacesstack{\mathcal{S}\!\mathit{paces}}
\def\Quotfunctor{\mathrm{Quot}}
\def\Hilbfunctor{\mathrm{Hilb}}
\def\Curvesstack{\mathcal{C}\!\mathit{urves}}
\def\Polarizedstack{\mathcal{P}\!\mathit{olarized}}
\def\Complexesstack{\mathcal{C}\!\mathit{omplexes}}
\def\Pic{\mathop{\mathrm{Pic}}\nolimits}
\def\Picardstack{\mathcal{P}\!\mathit{ic}}
\def\Picardfunctor{\mathrm{Pic}}
\def\Deformationcategory{\mathcal{D}\!\mathit{ef}}
\usepackage{cleveref}
\crefname{originallemma}{original lemma}{original lemmas}
\Crefname{originallemma}{Original lemma}{Original lemmas}
\makeatletter
\newlabel{algebra-lemma-ring-colimit-fp}{{00QN}{source}{}{sourcepointer-1}{}}
\newlabel{algebra-lemma-ring-colimit-fp@cref}{{[originallemma][0][]00QN}{[1][0][]source}{}{}{}}
\makeatother
\crefname{originallemma}{original lemma}{original lemmas}
\Crefname{originallemma}{Original lemma}{Original lemmas}
\makeatletter
\newlabel{categories-definition-representable-morphism}{{001X}{source}{}{sourcepointer-2}{}}
\newlabel{categories-definition-representable-morphism@cref}{{[originallemma][0][]001X}{[1][0][]source}{}{}{}}
\makeatother
\crefname{originallemma}{original lemma}{original lemmas}
\Crefname{originallemma}{Original lemma}{Original lemmas}
\makeatletter
\newlabel{constructions-lemma-relative-glueing}{{01LH}{source}{}{sourcepointer-3}{}}
\newlabel{constructions-lemma-relative-glueing@cref}{{[originallemma][0][]01LH}{[1][0][]source}{}{}{}}
\makeatother
\crefname{originallemma}{original lemma}{original lemmas}
\Crefname{originallemma}{Original lemma}{Original lemmas}
\makeatletter
\newlabel{lemma-canonical-extension}{{0GE3}{source}{}{sourcepointer-4}{}}
\newlabel{lemma-canonical-extension@cref}{{[originallemma][0][]0GE3}{[1][0][]source}{}{}{}}
\makeatother
\crefname{originallemma}{original lemma}{original lemmas}
\Crefname{originallemma}{Original lemma}{Original lemmas}
\makeatletter
\newlabel{lemma-representable-limit-preserving}{{0GE4}{source}{}{sourcepointer-5}{}}
\newlabel{lemma-representable-limit-preserving@cref}{{[originallemma][0][]0GE4}{[1][0][]source}{}{}{}}
\makeatother
\crefname{originallemma}{original lemma}{original lemmas}
\Crefname{originallemma}{Original lemma}{Original lemmas}
\makeatletter
\newlabel{limits-proposition-characterize-locally-finite-presentation}{{01ZC}{source}{}{sourcepointer-6}{}}
\newlabel{limits-proposition-characterize-locally-finite-presentation@cref}{{[originallemma][0][]01ZC}{[1][0][]source}{}{}{}}
\makeatother
\crefname{originallemma}{original lemma}{original lemmas}
\Crefname{originallemma}{Original lemma}{Original lemmas}
\makeatletter
\newlabel{limits-remark-limit-preserving}{{05LX}{source}{}{sourcepointer-7}{}}
\newlabel{limits-remark-limit-preserving@cref}{{[originallemma][0][]05LX}{[1][0][]source}{}{}{}}
\makeatother
\crefname{originallemma}{original lemma}{original lemmas}
\Crefname{originallemma}{Original lemma}{Original lemmas}
\makeatletter
\newlabel{spaces-limits-lemma-base-change-locally-finite-presentation}{{049M}{source}{}{sourcepointer-8}{}}
\newlabel{spaces-limits-lemma-base-change-locally-finite-presentation@cref}{{[originallemma][0][]049M}{[1][0][]source}{}{}{}}
\makeatother
\crefname{originallemma}{original lemma}{original lemmas}
\Crefname{originallemma}{Original lemma}{Original lemmas}
\makeatletter
\newlabel{spaces-limits-lemma-fibre-product-locally-finite-presentation}{{0GDZ}{source}{}{sourcepointer-9}{}}
\newlabel{spaces-limits-lemma-fibre-product-locally-finite-presentation@cref}{{[originallemma][0][]0GDZ}{[1][0][]source}{}{}{}}
\makeatother
\crefname{originallemma}{original lemma}{original lemmas}
\Crefname{originallemma}{Original lemma}{Original lemmas}
\makeatletter
\newlabel{spaces-limits-lemma-locally-finite-presentation-permanence}{{0GDY}{source}{}{sourcepointer-10}{}}
\newlabel{spaces-limits-lemma-locally-finite-presentation-permanence@cref}{{[originallemma][0][]0GDY}{[1][0][]source}{}{}{}}
\makeatother
\title{Representability of a transformation detected on Noetherian affines}
\author{The Stacks Project Authors}
\date{stacks-project:artin.tex:local-snapshot-20261012; GNU FDL 1.2 or later}
\begin{document}
\maketitle
\noindent\textit{Editorial note.} Statement span is lines 4886-4903 of sources/stacks-artin/pinned-original-source.tex; proof span is lines 4905-4993, the complete author proof (87 body lines).\par\medskip
\noindent\textit{Editorial note.} Navigation stubs: 10 cited-but-not-included labels are supplied as non-load-bearing pointers with their pinned tags where the snapshot tag index has them.\par\medskip
\noindent\textit{Editorial note.} extra\_packages is empty on purpose: the retained preamble already loads xy with options and a second declaration clashes in the assembler wrapper.\par\medskip
\noindent\textit{Editorial note.} Text is retained verbatim, including the author's own typos and placeholders.\par\medskip
\noindent\textit{Editorial note.} Queue index 4 of queue\_artin.json (pb=4309, stmt 4886-4903, proof 4905-4993).\par\medskip
\noindent\textit{Editorial note.} stacks-artin wave ar01 (2026-10-14): fresh adjudication accept per lane-s3/artin\_judgments.md; pool per Chief 10:5x direct dispatch (idle-first over S12's 01:5x note, SMC precedent); artin.tex 99 rows/67 with proof judged 46 accept / 21 reject / 5 AGP-overlap avoidance rows; patch-sd-envs contract spans scanned.\par\medskip
\noindent\textit{Allowed standard prerequisites (editorial).}\begin{itemize}
\item Explicit prerequisite boundary. The cited inputs are prerequisites; the argument in the retained span is not replaced by any cited result.
\end{itemize}\medskip
\ifdefined\hypertarget\hypertarget{sourcepointer-1}{}\fi
\noindent\textit{Source reference.} Original lemma 00QN: sources/stacks-artin/pinned-original-source.tex, source lines 4993--4993. Cited by the retained statement or proof but not part of the collected spans; supplied only as a navigation pointer. The argument does not depend on its content.\par\medskip
\ifdefined\hypertarget\hypertarget{sourcepointer-2}{}\fi
\noindent\textit{Source reference.} Original lemma 001X: sources/stacks-artin/pinned-original-source.tex, source lines 4993--4993. Cited by the retained statement or proof but not part of the collected spans; supplied only as a navigation pointer. The argument does not depend on its content.\par\medskip
\ifdefined\hypertarget\hypertarget{sourcepointer-3}{}\fi
\noindent\textit{Source reference.} Original lemma 01LH: sources/stacks-artin/pinned-original-source.tex, source lines 4993--4993. Cited by the retained statement or proof but not part of the collected spans; supplied only as a navigation pointer. The argument does not depend on its content.\par\medskip
\ifdefined\hypertarget\hypertarget{sourcepointer-4}{}\fi
\noindent\textit{Source reference.} Original lemma 0GE3: sources/stacks-artin/pinned-original-source.tex, source lines 4830--4840. Cited by the retained statement or proof but not part of the collected spans; supplied only as a navigation pointer. The argument does not depend on its content.\par\medskip
\ifdefined\hypertarget\hypertarget{sourcepointer-5}{}\fi
\noindent\textit{Source reference.} Original lemma 0GE4: sources/stacks-artin/pinned-original-source.tex, source lines 4858--4862. Cited by the retained statement or proof but not part of the collected spans; supplied only as a navigation pointer. The argument does not depend on its content.\par\medskip
\ifdefined\hypertarget\hypertarget{sourcepointer-6}{}\fi
\noindent\textit{Source reference.} Original lemma 01ZC: sources/stacks-artin/pinned-original-source.tex, source lines 4993--4993. Cited by the retained statement or proof but not part of the collected spans; supplied only as a navigation pointer. The argument does not depend on its content.\par\medskip
\ifdefined\hypertarget\hypertarget{sourcepointer-7}{}\fi
\noindent\textit{Source reference.} Original lemma 05LX: sources/stacks-artin/pinned-original-source.tex, source lines 4993--4993. Cited by the retained statement or proof but not part of the collected spans; supplied only as a navigation pointer. The argument does not depend on its content.\par\medskip
\ifdefined\hypertarget\hypertarget{sourcepointer-8}{}\fi
\noindent\textit{Source reference.} Original lemma 049M: sources/stacks-artin/pinned-original-source.tex, source lines 4993--4993. Cited by the retained statement or proof but not part of the collected spans; supplied only as a navigation pointer. The argument does not depend on its content.\par\medskip
\ifdefined\hypertarget\hypertarget{sourcepointer-9}{}\fi
\noindent\textit{Source reference.} Original lemma 0GDZ: sources/stacks-artin/pinned-original-source.tex, source lines 4993--4993. Cited by the retained statement or proof but not part of the collected spans; supplied only as a navigation pointer. The argument does not depend on its content.\par\medskip
\ifdefined\hypertarget\hypertarget{sourcepointer-10}{}\fi
\noindent\textit{Source reference.} Original lemma 0GDY: sources/stacks-artin/pinned-original-source.tex, source lines 4993--4993. Cited by the retained statement or proof but not part of the collected spans; supplied only as a navigation pointer. The argument does not depend on its content.\par\medskip
\begin{lemma}
\label{lemma-representable}
Let $\tau \in \{Zariski, \etale, smooth, syntomic, fppf\}$.
Let $F', G' : (\Sch/S)_\tau^{opp} \to \textit{Sets}$ be limit preserving
and sheaves. Let $a' : F' \to G'$ be a transformation of functors.
Denote $a : F \to G$ the restriction of $a' : F' \to G'$ to
$(\textit{Noetherian}/S)_\tau$. The following are equivalent
\begin{enumerate}
\item $a'$ is representable (as a transformation of functors, see
Categories, Definition \ref{categories-definition-representable-morphism}), and
\item for every object $V$ of $(\textit{Noetherian}/S)_\tau$
and every map $V \to G$ the fibre product
$F \times_G V : (\textit{Noetherian}/S)_\tau^{opp} \to \textit{Sets}$
is a representable functor, and
\item same as in (2) but only for $V$ affine finite type over $S$
mapping into an affine open of $S$.
\end{enumerate}
\end{lemma}

\begin{proof}
Assume (1). By Limits of Spaces, Lemma
\ref{spaces-limits-lemma-locally-finite-presentation-permanence}
the transformation $a'$ is limit preserving\footnote{This
makes sense even if $\tau \not = fppf$ as the underlying
category of $(\Sch/S)_\tau$ equals the underlying category
of $(\Sch/S)_{fppf}$ and the statement doesn't refer to the topology.}.
Take $\xi : V \to G$ as in (2). Denote $V' = V$ but viewed as an
object of $(\Sch/S)_\tau$. Since $G$ is the restriction of
$G'$ to $(\textit{Noetherian}/S)_\tau$ we see that
$\xi \in G(V)$ corresponds to $\xi' \in G'(V')$.
By assumption $V' \times_{\xi', G'} F'$ is representable
by a scheme $U'$. The morphism of schemes $U' \to V'$ corresponding to
the projection $V' \times_{\xi', G'} F' \to V'$ is locally of finite
presentation by
Limits of Spaces, Lemma
\ref{spaces-limits-lemma-base-change-locally-finite-presentation} and
Limits, Proposition
\ref{limits-proposition-characterize-locally-finite-presentation}.
Hence $U'$ is a locally Noetherian scheme and therefore $U'$ is
isomorphic to an object $U$ of $(\textit{Noetherian}/S)_\tau$.
Then $U$ represents $F \times_G V$ as desired.

\medskip\noindent
The implication (2) $\Rightarrow$ (3) is immediate. Assume (3).
We will prove (1). Let $T$ be an object of $(\Sch/S)_\tau$
and let $T \to G'$ be a morphism. We have to show
the functor $F' \times_{G'} T$ is representable by a scheme $X$
over $T$. Let $\mathcal{B}$ be the set of affine opens of
$T$ which map into an affine open of $S$. This is a basis
for the topology of $T$. Below we will show that for $W \in \mathcal{B}$
the fibre product $F' \times_{G'} W$ is representable
by a scheme $X_W$ over $W$. If $W_1 \subset W_2$ in $\mathcal{B}$, then
we obtain an isomorphism $X_{W_1} \to X_{W_2} \times_{W_2} W_1$ because both
$X_{W_1}$ and $X_{W_2} \times_{W_2} W_1$ represent the functor
$F' \times_{G'} W_1$.
These isomorphisms are canonical and satisfy the cocycle condition
mentioned in Constructions, Lemma
\ref{constructions-lemma-relative-glueing}.
Hence we can glue the schemes $X_W$ to a scheme $X$ over $T$.
Compatibility of the glueing maps with the maps
$X_W \to F'$ provide us with a map $X \to F'$.
The resulting map $X \to F' \times_{G'} T$ is an
isomorphism as we may check this locally on $T$ (as source
and target of this arrow are sheaves for the Zariski topology).

\medskip\noindent
Let $W$ be an affine scheme which maps into an affine open $U \subset S$.
Let $W \to G'$ be a map.
Still assuming (3) we have to show that $F' \times_{G'} W$
is representable by a scheme.
We may write $W = \lim V'_i$ as a directed limit of affine
schemes $V'_i$ of finite presentation over $U$, see
Algebra, Lemma \ref{algebra-lemma-ring-colimit-fp}.
Since $V'_i$ is of finite type over an Noetherian scheme,
we see that $V'_i$ is a Noetherian scheme.
Denote $V_i = V'_i$ but viewed as an object of
$(\textit{Noetherian}/S)_\tau$. As $G'$
is limit preserving can choose an $i$ and a map
$V'_i \to G'$ such that $W \to G'$ is the composition
$W \to V'_i \to G'$. Since $G$ is the restriction of $G'$
to $(\textit{Noetherian}/S)_\tau$ the morphism $V'_i \to G'$
is the same thing as a morphism $V_i \to G$ (see above).
By assumption (3) the functor $F \times_G V_i$ is representable by an object
$X_i$ of $(\textit{Noetherian}/S)_\tau$.
The functor $F \times_G V_i$ is limit preserving
as it is the restriction of $F' \times_{G'} V'_i$
and this functor is limit preserving by
Limits of Spaces, Lemma
\ref{spaces-limits-lemma-fibre-product-locally-finite-presentation},
the assumption that $F'$ and $G'$ are limit preserving, and
Limits, Remark \ref{limits-remark-limit-preserving} which
tells us that the functor of points of $V'_i$ is limit preserving.
By Lemma \ref{lemma-representable-limit-preserving}
we conclude that $X_i$ is locally of finite presentation over $S$.
Denote $X'_i = X_i$ but viewed as an object of
$(\Sch/S)_\tau$. Then we see that $F' \times_{G'} V'_i$
and the functors of points $h_{X'_i}$ are both extensions
of $h_{X_i} : (\textit{Noetherian}/S)_\tau^{opp} \to \textit{Sets}$
to limit preserving sheaves on $(\Sch/S)_\tau$.
By the equivalence of categories of Lemma \ref{lemma-canonical-extension}
we deduce that $X'_i$ represents $F' \times_{G'} V'_i$.
Then finally
$$
F' \times_{G'} W = F' \times_{G'} V'_i \times_{V'_i} W =
X'_i \times_{V'_i} W
$$
is representable as desired.
\end{proof}

\end{document}
